\section{Evidencia de warm start: cómo deben leerse las curvas}

El curso usa dos grupos de experimentos para mostrar que la inicialización puede ser útil entre tareas: las trayectorias de perplejidad de un modelo de lenguaje de una sola capa y la clasificación de MNIST con una y dos capas. Lo importante de las gráficas es el punto de partida y la trayectoria de optimización, no anunciar que la nueva arquitectura ya alcanzó la calidad de los modelos grandes.

\subsection{Modelo de lenguaje de una sola capa: el punto de partida es más bajo y la trayectoria sigue mejorando}

La sección anterior solo presentó el método para calcular los parámetros; esta sección pasa a la pregunta empírica más directa: ¿este valor inicial data-informed está realmente más cerca de un modelo utilizable que un punto de partida aleatorio? Al leer las curvas, primero se compara el punto de partida en la epoch 0 y después la trayectoria de descenso y la posición de early-stopping bajo la misma update rule; no se debe mirar primero solo el último punto e ignorar que las warm-start statistics también tienen un costo de cómputo.

La Perplexity (perplejidad) es la exponencial de la cross-entropy promedio:
\[
\operatorname{PPL}=\exp\!\left(-\frac{1}{T}\sum_{t=1}^{T}\log p(x_t\mid x_{<t})\right).
\]
Puede entenderse intuitivamente como el «número efectivo de candidatos» que el modelo enfrenta en cada token; mientras más baja, por lo general mejor. Sin embargo, la PPL obtenida con distinta tokenización, distintos datos o distinta configuración de context no puede compararse directamente.

\lecturefigure{slide-09-single-layer-warm-starts.jpg}{Trayectorias de entrenamiento con warm start y cold start en el modelo de una sola capa}{Grabación oficial de Stanford Online 00:24:03}

Al leer la figura, primero se observa la epoch 0: la PPL del cold start está cerca del vocabulary size, mientras que la del warm start ya es claramente menor. Después se observan la pendiente y el punto final: bajo la misma regla de iterative update, la curva de warm start sigue descendiendo y, al detenerse de forma temprana, mantiene una PPL más baja. Esto respalda que «un valor inicial explícito mejora el punto de partida y la optimización posterior», pero no controla todos los componentes de un Transformer moderno.

\begin{knowledgebox}{Lectura de la figura: esta gráfica no tiene self-attention}
En la sesión de Q\&A se aclaró que se trata de un feed-forward experiment con context one-hot concatenado. Las features de entrada ya están separadas, por lo que casi no se necesita attention para deshacer la superposition. Lo que valida es la single-layer explicit initialization, no una ablación completa de la attention de SAFFU.
\end{knowledgebox}

\begin{warningbox}{Usar la misma learning rate no equivale a una comparación completamente justa}
Warm start y random start pueden requerir distintos learning-rate schedule, regularization o paciencia de early-stopping. Una comparación rigurosa debe reportar al mismo tiempo una tuned cold baseline, varias seeds, el wall-clock y el costo de leer los datos y de calcular los estadísticos explícitos.
\end{warningbox}

\teachervoice{00:22:47--00:24:24, el ponente enfatiza que el warm start reduce la PPL desde antes del entrenamiento y que después sigue mejorando con la misma learning rate; si se usa early stopping, este también se activa antes y con una PPL más baja.}

\teachervoice{01:07:53--01:09:37, en la sesión de Q\&A se corrigió una lectura equivocada frecuente: esta gráfica warm/cold no tiene attention; el trabajo inicial era solo un feed-forward de una sola capa, y la solución explícita para la attention de SAFFU se agregó después.}

\subsection{MNIST: muestra que la fórmula a nivel de layer no se limita al texto}

Los pixels de MNIST son no negativos, así que el mismo warm start de la softmax layer puede aplicarse a image classification. En la figura, los modelos de una y dos capas con warm start alcanzan antes una accuracy alta; sin embargo, el modelo no incluye la convolution, la normalization ni el pipeline de augmentation de las redes de visión modernas.

\lecturefigure{slide-10-image-classification-warm-starts.jpg}{Curvas de clasificación con warm start en MNIST}{Grabación oficial de Stanford Online 00:30:45}

\begin{importantbox}{La conclusión mínima que respaldan los resultados de MNIST}
La softmax initialization explícita puede recibir features no lingüísticas y no negativas, y mejora el punto de partida en esta configuración de clasificación. No demuestra que la misma fórmula ya cubra convolution, ViT, diffusion ni datos visuales del mundo real.
\end{importantbox}

\teachervoice{00:29:14--00:31:20, el ponente usa MNIST para mostrar la layer-level generality de la fórmula; no afirma haber construido una vision architecture completa de alto rendimiento.}

\teachervoice{01:09:42--01:11:46, en la sesión de Q\&A se acotó todavía más el alcance: para manejar convolution u otras activation, todavía hay que derivar por separado el warm start de esas layer; de lo contrario, los parámetros aleatorios vuelven a introducir ruido.}

\subsection{Qué baselines deben agregarse al experimento}

En este tipo de investigación sobre initialization no basta con mirar la accuracy/PPL final. Hay que separar al menos cuatro tipos de costo: el cálculo de los estadísticos explícitos, el tuning de la baseline aleatoria, la backpropagation posterior y el efecto de fijar o actualizar los embedding sobre la caché.

\begin{knowledgebox}{Matriz experimental mínima}
\begin{tabularx}{\textwidth}{L{0.18\textwidth} L{0.24\textwidth} X}
\toprule
Comparación & Elementos fijos & Qué se debe reportar \\
\midrule
Cold vs warm & architecture, data, optimizer & seeds, loss inicial/final, wall-clock \\
Warm only vs warm+tune & initialization statistics & costo adicional de gradiente y ganancia en generalización \\
Frozen vs thawed embeddings & warm-start weights & aciertos de cache, velocidad de entrenamiento, calidad final \\
Local explicit vs end-to-end & parameter count & compositional gap y failure cases \\
\bottomrule
\end{tabularx}
\end{knowledgebox}

\subsection{Resumen de la sección}

Los experimentos de warm start muestran un mejor punto de partida y una mejor trayectoria de optimización, y en MNIST exhiben una capacidad limitada de transferirse a otro dominio. Lo que validan son condiciones específicas de layer y de datos, no que «todo el aprendizaje profundo pueda sustituirse por matrices de coocurrencia».
